\chapter{Muscle Twitching}

\section*{Definition}
Muscle twitching, or fasciculations, refers to spontaneous, involuntary contractions of small groups of muscle fibers visible beneath the skin as fine, flickering movements without joint movement. They result from spontaneous depolarization of motor units and are typically benign when isolated but may signify motor neuron pathology when widespread or persistent.

\section*{Pathophysiology}
Fasciculations arise from ectopic action potentials generated at any point along the lower motor neuron, from the anterior horn cell to the terminal axon branches. Increased neuronal excitability due to electrolyte imbalances, neurotransmitter alterations, or axonal irritation triggers spontaneous motor unit firing. Benign fasciculations involve hyperexcitable peripheral nerve terminals, while pathologic fasciculations reflect ongoing denervation and reinnervation in disorders affecting the anterior horn cell or axon.

\section*{Etiology}
\begin{itemize}
  \item Benign causes: Benign fasciculation syndrome, stress, anxiety, fatigue, exercise-induced, caffeine or stimulant excess, electrolyte disturbances (hypomagnesemia, hypokalemia, hypocalcemia)
  \item Medications: Diuretics, corticosteroids, beta-agonists, cholinesterase inhibitors, lithium
  \item Neuromuscular disorders: Amyotrophic lateral sclerosis (ALS), spinal muscular atrophy, post-polio syndrome, peripheral neuropathies, radiculopathies
  \item Metabolic conditions: Hyperthyroidism, hypoglycemia, hyperparathyroidism, vitamin D or B12 deficiency
  \item Infectious causes: Tetanus, rabies, neuroborreliosis, West Nile virus
\end{itemize}

\section*{Indications for Evaluation}
Seek care if:
\begin{itemize}
  \item Twitching accompanied by muscle weakness, atrophy, or loss of function
  \item Widespread or progressive fasciculations involving multiple body regions
  \item Associated sensory loss, cramps, or gait disturbance
  \item Onset after new medication or illness
  \item Bulbar involvement (tongue fasciculations, dysarthria, dysphagia)
\end{itemize}

\section*{Related Symptoms}
\begin{itemize}
  \item Muscle cramps
  \item Muscle weakness
  \item Muscle atrophy
  \item Myalgia
  \item Tremor
  \item Spasticity
  \item Fatigue
  \item Paresthesias
\end{itemize}
\textbf{Note:} Benign fasciculations are common, affect up to 70\% of the population, and are not associated with weakness or atrophy; the presence of weakness, atrophy, or reflex changes suggests neurogenic pathology requiring electromyography and neurology referral.

\section*{Self-Care / Home Management}
\begin{itemize}
  \item Reduce caffeine, nicotine, and stimulant intake as these increase neuronal excitability.
  \item Ensure adequate sleep, stress management, and hydration to minimize exacerbating factors.
  \item Correct electrolyte imbalances through dietary adjustments or supplementation under medical guidance.
  \item Gentle stretching and magnesium supplementation may provide symptomatic relief in benign cases.
\end{itemize}

\section*{Diagnostic Approach}
\begin{itemize}
  \item History characterizes onset, distribution, triggers, and associated neurologic symptoms (weakness, sensory loss, atrophy).
  \item Neurologic examination assesses muscle bulk, tone, strength, reflexes, coordination, and sensory function.
  \item Laboratory testing includes electrolytes (Ca, Mg, K, PO4), thyroid function, vitamin B12, and creatine kinase.
  \item Electrodiagnostic studies (EMG/NCS) distinguish benign fasciculations from neurogenic disease by evaluating for fibrillation potentials, positive sharp waves, and reinnervation changes.
\end{itemize}

\section*{Treatment / Management Overview}
\begin{itemize}
  \item Benign fasciculations require reassurance and avoidance of triggers; no pharmacotherapy is needed.
  \item Symptomatic treatment for bothersome fasciculations includes magnesium, carbamazepine, or gabapentin in selected cases.
  \item Underlying metabolic or electrolyte abnormalities are corrected when identified.
  \item Neurogenic causes are managed with disease-specific therapies (riluzole for ALS, immunosuppression for inflammatory neuropathies) and multidisciplinary supportive care.
\end{itemize}

\clearpage
